%% CivicTwin: A Counterfactual Graph Learning Framework for
%% Ex-Ante Urban Policy Comparison
%%
%% Compile from the REPOSITORY ROOT so that the \input paths below resolve:
%%     pdflatex paper/main.tex
%%     bibtex   main
%%     pdflatex paper/main.tex
%%     pdflatex paper/main.tex

\documentclass[journal]{IEEEtran}

\usepackage{amsmath}
\usepackage{amssymb}
\usepackage{amsfonts}
\usepackage{graphicx}
\usepackage{booktabs}
\usepackage{multirow}
\usepackage{algorithm}
\usepackage{algorithmic}
\usepackage{url}
\usepackage{balance}
\usepackage[hidelinks]{hyperref}

\graphicspath{{paper_assets/figures/}{../paper_assets/figures/}}

\newcommand{\CGCEN}{CG-CEN}
\newcommand{\HT}{\mathrm{H\!+\!T}}
\newcommand{\DPI}{\mathrm{DPI}}
\newcommand{\ITE}{\mathrm{ITE}}
\newcommand{\STE}{\mathrm{STE}}

\begin{document}

\title{CivicTwin: A Counterfactual Graph Learning Framework\\
for Ex-Ante Urban Policy Comparison}

\author{Abhay~Mahantesh~Zalaki,
        Ayush~Raj,
        and~Prabin~S~M%
  \thanks{Manuscript prepared 31 August 2026. Source code, synthetic data
  generator, benchmark harness and all figure/table assets are available at
  \url{https://github.com/AbhayZ1/CivicTwin}.}%
}

\markboth{Preprint --- CivicTwin: Counterfactual Graph Learning for Ex-Ante Urban Policy Comparison}%
{Zalaki \MakeLowercase{\textit{et al.}}: CivicTwin: Counterfactual Graph Learning for Ex-Ante Urban Policy Comparison}

\maketitle

\begin{abstract}
Urban policy interventions---upzoning, inclusionary housing mandates,
land-value capture---do not respect neighbourhood boundaries. A rezoning in one
census tract changes rents in its neighbours, violating the Stable Unit
Treatment Value Assumption (SUTVA) that underpins most quantitative policy
evaluation. We present CivicTwin, a Causal Graph Counterfactual Exposure
Network (\CGCEN{}) that treats ex-ante urban policy comparison as a causal
inference problem under partial interference rather than as a point-forecasting
problem. \CGCEN{} factorises each one-step-ahead prediction into a direct term
$\Phi(X_{i,t}, P_{i,t})$, which carries the Individual Treatment Effect, and a
neighbour-aggregated interference term $\sum_{j \in \mathcal{N}(i)} W_{ij}
\Psi(X_{j,t}, P_{j,t})$, which carries the Spatial Interference Effect. The
split is enforced architecturally: $\Phi$ reads a node-local temporal encoding
that never touches a message-passing operator, so it cannot absorb neighbour
information. We pair the model with two policy-facing indices---a Housing
$+$ Transportation overburden index and a level-aware Displacement Pressure
Index---and evaluate on a controlled synthetic city with a known
data-generating process, across ten random seeds and four policy scenarios.

Our central empirical finding is a boundary condition rather than a blanket
claim. \CGCEN{} attains a pooled MAE of $8.88 \pm 2.40$ against $11.87 \pm
3.12$ for a spatial lag model, $20.03 \pm 7.97$ for a non-graph temporal MLP,
and $27.90 \pm 1.27$ for naive persistence. However, its advantage over the
linear spatial baseline is statistically significant only under active policy
intervention ($p = 0.0015$--$0.0033$ across the three intervention scenarios)
and is \emph{not} significant in the untreated equilibrium ($p = 0.127$).
Topology ablations show the graph is load-bearing: removing edges costs
$+2.89$ MAE and randomly rewiring them costs $+4.88$ MAE, so a wrong topology
is more damaging than no topology at all. We further report a negative result
that shaped the design: under a generator whose spillover is a single permanent
shock, no model---ours included---derives any measurable benefit from the
graph, because permanent early spillover collapses into a static per-node
intercept. Graph learning earns its keep for urban policy only when the
spillover process is genuinely dynamic.
\end{abstract}

\begin{IEEEkeywords}
Causal inference, spatial interference, SUTVA, graph neural networks,
urban policy evaluation, counterfactual analysis, displacement pressure,
housing affordability.
\end{IEEEkeywords}

\IEEEpeerreviewmaketitle

\section{Introduction}
\label{sec:introduction}

\IEEEPARstart{U}{rban} policy is enacted on parcels and districts, but its
consequences travel. When a municipality upzones a corridor, relaxes parking
minimums, or captures land value to fund transit, the resulting change in rents,
land values and accessibility does not stop at the boundary of the treated
zone. It diffuses to adjacent neighbourhoods through housing search, commuting
patterns, developer expectations and the local price of land. This diffusion is
not a nuisance to be controlled away; for many of the questions planners
actually ask---\emph{who is displaced, and where?}---it is the phenomenon of
interest.

This poses a problem for quantitative evaluation. The standard potential
outcomes framework rests on the Stable Unit Treatment Value Assumption
(SUTVA)~\cite{rubin1980}, which requires that a unit's outcome depend only on
its own treatment assignment. Spatial spillover violates SUTVA by construction.
A tract adjacent to an upzoned corridor is, in a meaningful sense, partially
treated even though it received no intervention. Estimators that ignore this
attribute the neighbour's rent increase either to noise or, worse, to the
control condition, biasing the estimated effect of the policy toward zero.

\subsection{Ex-Ante Comparison, Not Point Prediction}

A second and less-discussed problem concerns the framing of the modelling task
itself. Much of the machine learning literature on urban and mobility data
frames the objective as forecasting: predict next month's rents, next hour's
traffic, next year's population. Spatio-temporal graph networks have been
notably successful at exactly this~\cite{li2018dcrnn,yu2018stgcn}. But a
planning agency deciding between an inclusionary housing mandate and a
land-value capture scheme is not asking what rents will be in 2031. It is
asking which of several counterfactual worlds is preferable, and for whom.
These are different questions with different evidentiary standards. A point
forecast can be accurate and still be useless for a policy choice if it cannot
be conditioned on an intervention that has not yet happened.

We therefore adopt an explicitly \emph{ex-ante comparative} scope rule. The
quantity of interest is not $\hat{Y}_{i,t+1}$ in absolute terms but the
contrast between $\hat{Y}_{i,t+1}$ under an intervention and
$\hat{Y}_{i,t+1}$ under the untreated counterfactual, together with the
corresponding contrast at each of unit $i$'s neighbours. Forecast accuracy
enters our evaluation not as the goal but as a necessary condition: a model
whose conditional expectations are badly calibrated cannot be trusted to
produce meaningful contrasts between counterfactuals.

\subsection{Approach}

We introduce the Causal Graph Counterfactual Exposure Network (\CGCEN{}), which
makes the interference structure explicit in the architecture rather than
implicit in a learned representation. Following the partial interference
formulation of Hudgens and Halloran~\cite{hudgens2008}, we assume a unit's
outcome depends on its own treatment and on a summary of its neighbours'
treatments mediated by the graph. \CGCEN{} factorises the one-step-ahead
prediction as
\begin{equation}
  \hat{Y}_{i,t+1} = \underbrace{\Phi(X_{i,t}, P_{i,t})}_{\text{direct}}
  + \underbrace{\sum_{j \in \mathcal{N}(i)} W_{ij} \cdot
  \Psi(X_{j,t}, P_{j,t})}_{\text{spatial interference}} ,
  \label{eq:decomposition-intro}
\end{equation}
where $\Phi$ estimates the Individual Treatment Effect (ITE) and $\Psi$,
aggregated over neighbours with row-normalised weights $W_{ij}$, captures the
Spatial Interference Effect (STE). The separation is enforced structurally:
$\Phi$ consumes a node-local temporal encoding that never passes through a
message-passing operator, so by construction it cannot launder neighbour
information into the direct term. Both effects are recovered by contrasting
each branch against its own $P = 0$ counterfactual, which makes the ITE exactly
zero on untreated units and the STE exactly zero on an edgeless graph---two
identities we verify numerically as regression tests.

Because credible spatial policy data with known ground truth is scarce, we
evaluate on a transparent synthetic city whose data-generating process we
control. This is a deliberate methodological choice. It lets us ask a question
that observational data cannot answer: \emph{when the spillover is known to
exist, does the graph model actually find it?} The answer turns out to be
conditional, and that condition is our main contribution.

\subsection{Related Work}

Our work sits at the intersection of three literatures.

\emph{Causal inference under interference} relaxes SUTVA by positing that a
unit's outcome depends on others' assignments. Hudgens and
Halloran~\cite{hudgens2008} formalise partial interference within disjoint
clusters; Ogburn and VanderWeele~\cite{ogburn2014} give a graphical account;
and Forastiere et al.~\cite{forastiere2021} extend identification to
observational network settings with an explicit exposure mapping. We adopt
their exposure-mapping formulation but embed it in a learned spatio-temporal
model rather than a parametric outcome regression.

\emph{Neural counterfactual estimation} learns representations for treatment
effect estimation, typically under SUTVA. Shalit et al.~\cite{shalit2017cfr}
give generalisation bounds for individual treatment effects, and Wager and
Athey~\cite{wager2018causalforest} provide asymptotic inference for forest-based
heterogeneous effects. Ma and Tresp~\cite{ma2021netdeconf} extend neural
counterfactual estimation to networked interference, which is closest to our
setting; our contribution relative to that line is the architectural
identification of the direct/interference split and the explicit temporal
dimension.

\emph{Spatio-temporal graph learning} has been driven largely by traffic
forecasting. Kipf and Welling~\cite{kipf2017gcn} introduced the graph
convolution we take as a reference point, and DCRNN~\cite{li2018dcrnn} and
STGCN~\cite{yu2018stgcn} established the recurrent and convolutional
spatio-temporal templates. These are point-forecasting models evaluated on
observational series; neither exposes a treatment channel or a counterfactual
contrast, which is what the policy-comparison setting requires.

On the measurement side, the combined housing-and-transportation affordability
framing follows the H$+$T Index~\cite{cnt2023ht}, and our displacement
indicators are informed by the empirical displacement literature, particularly
Zuk et al.~\cite{zuk2018gentrification}, who document how public investment can
simultaneously raise accessibility and displacement risk---the tension our
$\HT$ index is designed to make visible.

\subsection{Contributions}

\begin{enumerate}
  \item \textbf{An architecturally identified interference split.} We give a
  spatio-temporal architecture in which the direct and interference channels
  are separated by construction rather than by regularisation, together with
  counterfactual contrasts that recover ITE and STE with exact degenerate-case
  behaviour (Section~\ref{sec:methodology}).

  \item \textbf{A level-aware displacement index.} We show that a
  growth-rate-only Displacement Pressure Index is structurally incapable of
  registering a sustained policy effect: a permanent level shift produces a
  single transient spike and then returns to zero. We add a cumulative
  rent-burden term that restores sensitivity to persistent burden, increasing
  measured policy contrasts by more than two orders of magnitude
  (Sections~\ref{sec:methodology} and~\ref{sec:experiments}).

  \item \textbf{A conditional result on when graphs help.} \CGCEN{} improves
  significantly over a linear spatial lag model under active policy
  intervention pulses ($p = 0.0015$--$0.0033$), but \emph{not} in the untreated
  equilibrium ($p = 0.127$). We report this boundary explicitly rather than
  pooling it away, because the pooled comparison would overstate the
  method's generality.

  \item \textbf{A negative result with a mechanism.} Under a generator in which
  spillover arrives as a single permanent shock, no model we tested---including
  ours, and including the spatial lag model whose learned spatial coefficient
  converged to $\rho \approx 0$---derived measurable benefit from the graph.
  Removing edges \emph{improved} accuracy. We diagnose the cause: permanent
  early spillover is absorbed into each node's own level and becomes a static
  intercept, indistinguishable from node-specific effects. We show that making
  the shock process stochastic and recurrent restores identifiability
  (Section~\ref{sec:data-engine}).

  \item \textbf{A reproducible benchmark.} We release the generator, a
  four-model baseline suite, a ten-seed harness with confidence intervals and
  paired significance tests, and topology ablations (real, edgeless, randomly
  rewired), together with all figure and table assets used in this manuscript.
\end{enumerate}

The fourth contribution deserves emphasis. It is the kind of finding that is
usually discarded during model development, and discarding it here would have
produced a paper claiming that a graph neural network captures urban spillover
when in fact the benchmark contained no dynamic spillover for it to capture.
Reporting the boundary between the regimes is, we argue, more useful to the
field than reporting only the regime in which the method wins.

\subsection{Organisation}

Section~\ref{sec:methodology} formalises the graph, the causal decomposition
and the two policy indices. Section~\ref{sec:data-engine} describes the
synthetic data engine, the dynamic shock process and the logistic stabilisation
that keeps long-horizon trajectories bounded.
Section~\ref{sec:experiments} reports the ten-seed benchmark, significance
tests and topology ablations. Section~\ref{sec:discussion} discusses
limitations---chiefly the absence of a full-scale empirical
evaluation---and outlines the path to external validation.

\section{Formal Methodology and Causal Engine}
\label{sec:methodology}

\subsection{The Spatio-Temporal Graph}

Let $G = (V, E, W)$ be an undirected weighted graph in which each vertex
$i \in V$, $|V| = N$, is a neighbourhood unit (a census tract in the empirical
setting, a lattice cell in the synthetic setting) and each edge
$(i,j) \in E$ encodes spatial adjacency. We write $\mathcal{N}(i) =
\{ j : (i,j) \in E \}$ for the neighbour set of $i$ and $d_i = |\mathcal{N}(i)|$
for its degree. The weight matrix $W \in \mathbb{R}^{N \times N}$ is the
row-normalised adjacency,
\begin{equation}
  W_{ij} =
  \begin{cases}
    a_{ij} \big/ \sum_{k} a_{ik}, & (i,j) \in E,\\[2pt]
    0, & \text{otherwise,}
  \end{cases}
  \label{eq:weights}
\end{equation}
with $a_{ij}$ the raw adjacency entry, so that $\sum_{j} W_{ij} = 1$ for every
non-isolated $i$. Row normalisation makes the interference term an average
rather than a sum, which keeps its scale invariant to degree and prevents
high-degree units from dominating the objective.

Each unit carries a feature trajectory $X_{i,t} \in \mathbb{R}^{F}$ over
discrete steps $t = 0, \dots, T$. In this work $F = 6$: land value, rent,
population, median income, housing units and accessibility. A policy assignment
$P_{i,t} \in \{0,1\}$ indicates whether unit $i$ is subject to the intervention
at step $t$. We write $X_{i,t-W_h+1:t}$ for the length-$W_h$ history window
ending at $t$.

\subsection{Interference and the Exposure Mapping}

Under SUTVA, the potential outcome of unit $i$ would be written
$Y_i(P_i)$, depending on $i$'s assignment alone. Spatial spillover invalidates
this. We adopt instead the partial interference
assumption~\cite{hudgens2008,ogburn2014}: unit $i$'s outcome depends on its own
assignment and on its neighbours' assignments only through a low-dimensional
exposure summary. We use the degree-normalised exposure
\begin{equation}
  e_i = \sum_{j \in \mathcal{N}(i)} W_{ij} P_{j,t}
      = \frac{1}{d_i} \sum_{j \in \mathcal{N}(i)} P_{j,t},
  \label{eq:exposure}
\end{equation}
so that $e_i \in [0,1]$ measures the treated fraction of $i$'s neighbourhood.
The potential outcome becomes $Y_i(P_i, e_i)$, and the estimands of interest
separate into a direct effect holding exposure fixed and an interference effect
holding own-treatment fixed~\cite{forastiere2021}.

\subsection{Causal Effect Decomposition}

The \CGCEN{} forward pass realises this separation. Let $h_i \in
\mathbb{R}^{H}$ be a node-local temporal encoding of the history window,
produced by a shared feature encoder followed by a gated recurrent
unit~\cite{cho2014gru}:
\begin{align}
  z_{i,k} &= \mathrm{ReLU}\!\left( \mathbf{W}_{\mathrm{enc}} \tilde{X}_{i,k}
             + b_{\mathrm{enc}} \right),
             \quad k = t - W_h + 1, \dots, t, \label{eq:encoder}\\
  h_i     &= \mathrm{GRU}\!\left( z_{i,t-W_h+1}, \dots, z_{i,t} \right),
             \label{eq:gru}
\end{align}
where $\tilde{X}$ denotes the $z$-score standardised features of
Section~\ref{subsec:scaling}. Critically, \eqref{eq:encoder}--\eqref{eq:gru}
contain no message passing: $h_i$ is a function of unit $i$'s own history only.

Writing $s_i = [\, h_i \,\Vert\, P_{i,t} \,]$ for the state-treatment
concatenation, the prediction is
\begin{equation}
  \hat{Y}_{i,t+1}
  = \underbrace{X_{i,t} + \mathrm{MLP}_{\Phi}(s_i)}_{\Phi(X_{i,t}, P_{i,t})}
  + \underbrace{\sum_{j \in \mathcal{N}(i)} W_{ij} \,
    \mathrm{MLP}_{\Psi}(s_j)}_{\sum_j W_{ij} \Psi(X_{j,t}, P_{j,t})} .
  \label{eq:decomposition}
\end{equation}
The term $X_{i,t}$ inside $\Phi$ is a residual anchor: the network predicts the
\emph{change} from the last observation rather than the level. This is not a
cosmetic choice. Without it the GRU must reconstruct an absolute level through a
saturating $\tanh$ bottleneck. Ablating the anchor while holding everything else
fixed degrades pooled MAE from $8.92$ to $37.28$, a factor of $4.2$; the
near-linear ``next $\approx$ last'' signal that dominates these series is
destroyed by the recurrent nonlinearity. Because
$X_{i,t}$ is a function of unit $i$'s own state, the anchor is properly part of
the direct term and does not contaminate the interference channel.

The identification of the split is architectural. $\mathrm{MLP}_{\Phi}$ reads
only $s_i$, which derives from the message-passing-free encoding
\eqref{eq:gru}; it therefore has no path by which neighbour state could enter.
$\mathrm{MLP}_{\Psi}$ is evaluated at the neighbours and aggregated with $W$;
it is the only route by which neighbour information reaches the prediction.

\subsection{Recovering ITE and STE}

Effects are recovered by contrasting each branch against its own untreated
counterfactual. Let $s_i^{0} = [\, h_i \,\Vert\, 0 \,]$. Then
\begin{align}
  \tau_i &= \mathrm{MLP}_{\Phi}(s_i) - \mathrm{MLP}_{\Phi}(s_i^{0}),
  \label{eq:ite}\\
  \eta_i &= \sum_{j \in \mathcal{N}(i)} W_{ij}
            \left( \mathrm{MLP}_{\Psi}(s_j) - \mathrm{MLP}_{\Psi}(s_j^{0}) \right),
  \label{eq:ste}
\end{align}
where $\tau_i$ is the estimated ITE and $\eta_i$ the estimated STE. Two
identities follow immediately and serve as correctness checks in our test
suite. First, if $P_{i,t} = 0$ then $s_i = s_i^{0}$ and $\tau_i = 0$ exactly:
an untreated unit has no direct effect. Second, if $E = \varnothing$ then the
aggregation in \eqref{eq:ste} is empty and $\eta_i = 0$ exactly: with no graph
there is no interference. The additive identity
$\hat{Y}_{i,t+1} = \Phi_i + \sum_j W_{ij}\Psi_j$ holds by construction to
floating-point tolerance.

Algorithm~\ref{alg:causal_stgnn} states the full forward pass together with the
counterfactual contrasts.

\begin{algorithm}[t]
\caption{Causal ST-GNN forward pass with direct/spillover decomposition (CG-CEN)}
\label{alg:causal_stgnn}
\begin{algorithmic}[1]
\REQUIRE Node history $X_{i,t-W+1:t} \in \mathbb{R}^{W \times F}$ for all $i \in \mathcal{V}$;
         policy assignment $P_{i,t} \in \{0,1\}$; adjacency $\mathcal{N}(\cdot)$ with
         row-normalised weights $W_{ij}$
\ENSURE  Prediction $\hat{Y}_{i,t+1}$, individual treatment effect $\tau_i$,
         spatial interference effect $\eta_i$
\STATE $\mu, \sigma \gets$ \textsc{FitScaler}$(X_{\text{train}})$ \COMMENT{training split only}
\STATE $Z_{i,\cdot} \gets (X_{i,\cdot} - \mu) / \sigma$ for all $i$
\FOR{$k = t-W+1$ \TO $t$}
    \STATE $e_{i,k} \gets \mathrm{ReLU}(\mathbf{W}_{\text{enc}} Z_{i,k} + b_{\text{enc}})$
\ENDFOR
\STATE $h_i \gets \mathrm{GRU}\big(e_{i,t-W+1}, \dots, e_{i,t}\big)$ \COMMENT{node-local; no message passing}
\STATE $s_i \gets [\,h_i \,\Vert\, P_{i,t}\,]$
\STATE $\Phi_i \gets \mathrm{MLP}_{\Phi}(s_i) + X_{i,t}$ \COMMENT{direct term, residual anchor}
\STATE $\Psi_j \gets \mathrm{MLP}_{\Psi}(s_j)$ for all $j$
\STATE $\Omega_i \gets \sum_{j \in \mathcal{N}(i)} W_{ij} \, \Psi_j$ \COMMENT{spatial interference}
\STATE $\hat{Y}_{i,t+1} \gets \Phi_i + \Omega_i$
\STATE $s_i^{0} \gets [\,h_i \,\Vert\, 0\,]$ \COMMENT{counterfactual: no intervention}
\STATE $\tau_i \gets \mathrm{MLP}_{\Phi}(s_i) - \mathrm{MLP}_{\Phi}(s_i^{0})$ \COMMENT{ITE}
\STATE $\eta_i \gets \sum_{j \in \mathcal{N}(i)} W_{ij}\big(\mathrm{MLP}_{\Psi}(s_j)
        - \mathrm{MLP}_{\Psi}(s_j^{0})\big)$ \COMMENT{STE}
\RETURN $\hat{Y}_{i,t+1}, \tau_i, \eta_i$
\end{algorithmic}
\end{algorithm}


\subsection{Feature Scaling}
\label{subsec:scaling}

The panel spans several orders of magnitude---median income is $O(10^{4})$
while accessibility is $O(10^{0})$---so an unscaled squared-error objective is
dominated by the large-magnitude columns and the gradient signal for the
remainder is effectively lost. We apply per-feature $z$-score standardisation,
$\tilde{X}_{\cdot,\cdot,f} = (X_{\cdot,\cdot,f} - \mu_f)/\sigma_f$, with
$\mu_f$ and $\sigma_f$ estimated \emph{on the training split only} and reused
unchanged for the test split, so that no distributional information leaks
backwards across the temporal boundary. Constant features are assigned
$\sigma_f = 1$ to avoid division by zero. All reported metrics are computed
after inverting the transform, so errors are in the original units and remain
comparable across configurations.

\subsection{Baseline Models}

We compare against three baselines spanning the relevant hypothesis space.
The \emph{Temporal MLP} flattens the history window and maps it to the next
step through a two-layer perceptron; it is graph-blind by construction and
isolates the value of the topology. The \emph{Spatial Lag Model} (SLM) is a
learned analogue of the classical spatial autoregressive specification of
spatial econometrics~\cite{anselin1988,lesage2009}, predicting
\begin{equation}
  \hat{Y}_{i,t+1} = X_{i,t} + \beta^{\top} \mathrm{vec}(X_{i,t-W_h+1:t})
  + \rho \odot \sum_{j \in \mathcal{N}(i)} W_{ij} X_{j,t},
  \label{eq:slm}
\end{equation}
with a learned spatial coefficient $\rho$ and a residual anchor matching
\eqref{eq:decomposition} so that the comparison isolates functional form rather
than parameterisation. The \emph{Naive Persistence} model predicts
$\hat{Y}_{i,t+1} = X_{i,t}$; it is parameter-free and never trained. It is the
floor every learned model must clear before any claim about spillover structure
is meaningful, and including it is what revealed, in an earlier iteration of
this work, that two learned models were both losing to it.

\subsection{Policy Indices}

\subsubsection{Housing $+$ Transportation Overburden Index}

Housing cost alone understates the burden borne by households in
poorly-connected neighbourhoods, because low rents are frequently purchased
with high commuting costs. Following the combined housing-and-transportation
affordability convention, we define
\begin{equation}
  \HT_{i,t}
  = \frac{\mathrm{Rent}_{i,t} + \mathrm{TransitCost}_{i,t}}
         {\mathrm{MedianIncome}_{i,t}} ,
  \label{eq:ht}
\end{equation}
where transit cost falls with accessibility,
\begin{equation}
  \mathrm{TransitCost}_{i,t} =
  \frac{c_0}{\max(\mathrm{Accessibility}_{i,t}, \epsilon)^{\kappa}} ,
  \label{eq:transit}
\end{equation}
with base cost $c_0$ and elasticity $\kappa$. This is precisely the mechanism a
rent-to-income ratio misses: a transit investment that raises accessibility
lowers $\HT$ even when it raises rent, and whether the net effect is
progressive or regressive is an empirical question the index can answer.

\subsubsection{Level-Aware Displacement Pressure Index}

The Displacement Pressure Index blends adverse movements in rent, low-income
share and accessibility:
\begin{align}
  \DPI_{i,t} =\;& \alpha \left( \frac{\Delta \mathrm{Rent}_{i,t}}
                                    {\mathrm{Rent}_{i,t-1}} \right)
  + \beta \left( 1 - \frac{\mathrm{LowIncomeShare}_{i,t}}
                          {\mathrm{LowIncomeShare}_{i,t-1}} \right) \nonumber\\
  &+ \gamma \left( \Delta \mathrm{Accessibility}_{i,t} \right)
   + \delta \, B_{i,t} .
  \label{eq:dpi}
\end{align}
The first three terms are the growth-rate formulation. The fourth is a
correction we introduce. A growth-rate index is, by construction, blind to
sustained level differences: a policy that permanently raises rents by six per
cent changes the growth rate only in the single period in which the shift
occurs, after which the index returns to its baseline path. In our experiments
this rendered every policy contrast on the order of $10^{-4}$---numerically
present but substantively meaningless. We therefore add a cumulative
rent-burden level term measured against each unit's own initial burden,
\begin{equation}
  B_{i,t} = \frac{\HT_{i,t}}{\max(\HT_{i,0}, \epsilon)} - 1 ,
  \label{eq:burden}
\end{equation}
which is zero at $t = 0$ by construction and retains a permanent offset under a
permanent shift. Setting $\delta = 0$ recovers the pure growth-rate
specification exactly, a property we pin with a regression test so that the
original formulation remains verifiable.

Both indices are constructed comparative indicators over a controlled panel.
Neither is a calibrated forecast of realised displacement, and we use them only
for ranking scenarios against a common baseline.

\section{Synthetic Data Engine and Dynamic Equilibrium}
\label{sec:data-engine}

\subsection{Rationale for a Controlled Generator}

Evaluating a spillover-aware estimator on observational urban data is
circular: the spillover structure is exactly what is unknown, so a model that
appears to exploit it cannot be distinguished from one that has fit
confounding. We therefore evaluate primarily on a synthetic city whose
data-generating process we specify completely. This licenses the question that
motivates the paper---\emph{when spillover is known to be present, does the
graph model recover it?}---and, as Section~\ref{subsec:collapse} shows, the
answer is not automatically yes.

\subsection{Graph and State Variables}

The synthetic city is a $4 \times 4$ lattice of $N = 16$ neighbourhood units
with four-connected adjacency, giving $|E| = 48$ directed edge entries and
degrees between two and four. Each unit carries the six state variables of
Section~\ref{sec:methodology}, initialised with multiplicative Gaussian
heterogeneity around city-wide bases (land value $100$, rent $50$, population
$2000$, median income $60{,}000$, housing units $900$, accessibility $1.0$).
Generation is fully seeded and reproducible.

Dynamics are deliberately transparent. Accessibility shocks raise land values;
land value growth feeds rent growth with a one-step lag; rent pressure damps
population, income and housing growth; and neighbour terms diffuse land and
rent across edges. A designated shock set $S = \{0,1,4,7\}$ receives a permanent
accessibility uplift from $t \geq 3$, representing a transit investment.

\subsection{The Static Intercept Collapse}
\label{subsec:collapse}

Our first generator used only that permanent shock. Under it, every model we
tested failed to benefit from the graph. The Spatial Lag Model's learned
spatial coefficient converged to $\rho \approx 0$ across seeds
($-0.05$, $-0.07$, $0.02$); \CGCEN{} with edges removed was \emph{better} than
with the real graph ($14.16$ versus $14.82$ MAE); and randomly rewiring the
edges was statistically indistinguishable from the true topology
($p = 0.70$--$0.93$).

The mechanism is instructive. A permanent shock applied at $t = 3$ to a fixed
node set means that, by the evaluation window at $t \geq 12$, each unit's
neighbour contribution has converged to a time-invariant constant. That
constant is absorbed into the unit's own level and is recovered perfectly by
the residual anchor $X_{i,t}$ in \eqref{eq:decomposition}. In other words,
spillover that is both permanent and early is observationally equivalent to a
node-specific fixed effect. There is nothing for a graph to contribute, because
the information has already been encoded in the ego trajectory. We call this
\emph{static intercept collapse}. It is, we suspect, a common and under-reported
failure mode in spatio-temporal benchmarks that inject spillover as a one-time
regime change.

\subsection{Dynamic Network Shocks}

The remedy is to make the spillover process genuinely dynamic. We augment the
generator with stochastic, recurrent accessibility pulses. At each step every
unit independently receives a pulse with probability $p_{\mathrm{pulse}} =
0.22$, of magnitude drawn uniformly from $[0.5\,\sigma_p,\, 1.5\,\sigma_p]$
with $\sigma_p = 1.2$:
\begin{equation}
  u_{i,t} \sim \mathrm{Bernoulli}(p_{\mathrm{pulse}}) \cdot
                \mathcal{U}(0.5\sigma_p,\, 1.5\sigma_p).
  \label{eq:pulse}
\end{equation}
The pulse enters the unit's own observed accessibility at $t$, and---this is
the essential coupling---transmits to \emph{neighbours'} rent growth at $t+1$:
\begin{equation}
  g^{\mathrm{rent}}_{i,t} \mathrel{+}= \lambda_p \sum_{j \in \mathcal{N}(i)}
  W_{ij} \, u_{j,t-1},
  \label{eq:pulse-transmission}
\end{equation}
with transmission strength $\lambda_p = 0.22$.

Because the pulse is transient and stochastic, a neighbour's pulse at $t$
cannot be inferred from unit $i$'s own history: it is new information available
only across an edge. This is what makes the graph necessary rather than merely
available. Accessibility is an observed feature, so the neighbour's pulse is
visible to a message-passing model at $t$ and predictive of ego rent at $t+1$,
while remaining invisible to any purely node-local estimator.

\subsection{Logistic Carrying Capacities and Bounded Trajectories}

The original generator applied purely multiplicative growth with no
equilibrating force, which produced a second and more damaging defect. Rents
compounded without bound; as the rent-to-income ratio grew, the income update's
rent-pressure term drove income sharply negative; and income collided with its
hard floor. In the unmodified generator, income reached the floor at $t = 52$
and by $t = 100$ \emph{all sixteen units were pinned at it} while mean rent had
compounded to $\approx 4 \times 10^{10}$. Any statistic computed beyond
$t \approx 52$ measured the clamp function, not the economy.

We introduce a logistic carrying-capacity structure. Each growth rate is
multiplied by a headroom factor that vanishes as the state approaches its
capacity,
\begin{equation}
  \mathrm{headroom}(x, K) = \max\!\left(0,\; 1 - \frac{x}{K}\right),
  \label{eq:headroom}
\end{equation}
with land capacity $K_L = 12 \times$ base land value, income capacity
$K_Y = 4 \times$ base income, and a rent capacity tied to contemporaneous
ability to pay, $K_R = 0.012 \times \mathrm{Income}_{i,t-1}$. Rent and income
additionally carry mean-reversion terms pulling them toward
$0.55\,K_R$ and toward base income respectively, with rates $0.05$ and $0.02$.

The economic interpretation is a supply response: as rents approach the local
affordability ceiling, further growth is choked off by construction, vacancy
and out-migration. The numerical effect is that the synthetic economy reaches a
dynamic equilibrium rather than diverging. Under the revised generator
\emph{no unit reaches any hard state floor through $t = 100$}: mean rent
settles at $\approx 1.6 \times 10^{3}$, income at $\approx 1.5 \times 10^{5}$
and $\HT$ at $\approx 0.011$. Figure~\ref{fig:trajectories} shows the resulting
trajectories.

\subsection{Policy Scenarios}

We define four scenarios, applied to the shock-node boundary $S$ as an
intervention on the generated panel. All are deliberately simple and additive
so that their ground-truth direction is unambiguous.

\begin{itemize}
  \item \textbf{Baseline.} No intervention. Provides the untreated
  counterfactual against which all contrasts are computed.

  \item \textbf{Market-Led Development.} Upzoning without affordability
  conditions: land values rise by $4\%$ and rents by $6\%$ per unit intensity
  from $t \geq 2$. Expected to \emph{increase} displacement pressure.

  \item \textbf{Inclusionary Housing.} A mandate coupling new development to
  affordable set-asides: rents fall by $6\%$ and housing stock rises by $8\%$
  per unit intensity from $t \geq 1$. Expected to \emph{reduce} pressure.

  \item \textbf{Land-Value Capture.} Taxation of land-value uplift recycled
  into affordability: land values rise by a damped $2\%$ while rents fall by
  $5\%$ per unit intensity from $t \geq 2$. Expected to reduce pressure, with a
  different land/rent signature from inclusionary housing.
\end{itemize}

Scenario intensities ($0.8$, $0.6$, $0.7$ respectively) are configuration
parameters, not fitted quantities. For each of ten seeds we export all four
scenario panels together with adjacency, edge index, low-income share and the
treated boundary, giving $10 \times 4 = 40$ reproducible datasets.

\begin{figure*}[t]
  \centering
  \includegraphics[width=\textwidth]{scenario_trajectories.pdf}
  \caption{Scenario trajectories over $t = 0 \dots 100$ (city-mean, seed 1).
  Panels (a) and (c) show $\HT$ and $\DPI$ levels; panels (b) and (d) show each
  intervention relative to the untreated baseline, which is the comparative
  quantity of interest. Under the logistic carrying capacities and mean
  reversion of \eqref{eq:headroom}, the economy reaches a dynamic equilibrium
  and no unit reaches a hard state floor over the full horizon. Policy effects
  persist as sustained level offsets rather than the single-step transients
  produced by a growth-rate-only index.}
  \label{fig:trajectories}
\end{figure*}

\section{Experimental Evaluation and Results}
\label{sec:experiments}

\subsection{Protocol}

Each run generates a seeded synthetic city of $N = 16$ units over $T = 18$
steps, applies one of the four policy scenarios, and forms rolling history
windows of length $W_h = 4$ with next-step targets. The split is strictly
temporal---the earliest $70\%$ of windows train, the remainder test---with no
shuffling, so no future information reaches the training set. The scaler of
Section~\ref{subsec:scaling} is fit on the training windows alone.

We restrict the modelled feature set to the three spillover-bearing channels
(land value, rent, accessibility). All four models are constructed under an
identical random seed immediately before instantiation, so that comparisons
between architectures, and between topology conditions, reflect the quantity
under study rather than the draw of the initial weights. Learned models train
for $250$ epochs with Adam at learning rate $10^{-3}$ under a mean squared
error objective; persistence is parameter-free and is never trained. Metrics
are computed after inverting the scaler, pooled over every window, unit and
feature, so the reported RMSE is a true root-mean-squared error rather than a
root-mean-square of per-window absolute errors.

The full design is $10$ seeds $\times$ $4$ scenarios $\times$ $4$ models for
the main sweep, plus $10$ seeds $\times$ $4$ scenarios $\times$ $2$ ablated
topologies for \CGCEN{}, giving $240$ fitted models. Aggregate statistics use
the sample standard deviation ($\mathrm{ddof} = 1$) and Student-$t$ $95\%$
confidence intervals. Model comparisons use two-sided paired $t$-tests over
seeds, pairing on the seed so that between-seed variance in task difficulty is
differenced out.

\subsection{Forecast Accuracy}

Table~\ref{tab:benchmark_results} reports per-scenario MAE and RMSE with
significance markers; Figure~\ref{fig:model-comparison} summarises the pooled
comparison.

\begin{table*}[t]
\centering
\caption{Forecast accuracy across policy scenarios on the synthetic CivicTwin panel. Values are mean $\pm$ standard deviation of the pooled test-set error over independent seeds. Significance markers denote a paired $t$-test of each baseline against the Causal ST-GNN on MAE ($^{*}p<0.05$, $^{**}p<0.01$, $^{***}p<0.001$).}
\label{tab:benchmark_results}
\small
\begin{tabular}{lcccccccc}
\toprule
\textbf{Model} & \multicolumn{2}{c}{\textbf{Baseline}} & \multicolumn{2}{c}{\textbf{Market-Led Development}} & \multicolumn{2}{c}{\textbf{Inclusionary Housing}} & \multicolumn{2}{c}{\textbf{Land-Value Capture}} \\
\cmidrule(lr){2-3}\cmidrule(lr){4-5}\cmidrule(lr){6-7}\cmidrule(lr){8-9}
 & MAE & RMSE & MAE & RMSE & MAE & RMSE & MAE & RMSE \\
\midrule
Causal ST-GNN (CG-CEN) & $9.901 \pm 2.234$ & $15.482 \pm 2.994$ & $8.671 \pm 2.954$ & $13.763 \pm 4.719$ & $8.645 \pm 2.368$ & $13.316 \pm 3.630$ & $8.303 \pm 1.991$ & $12.688 \pm 2.981$ \\
Spatial Lag Model & $11.837 \pm 3.234$ & $19.269 \pm 5.083$ & $12.014 \pm 3.283^{**}$ & $19.595 \pm 5.167$ & $11.796 \pm 3.231^{**}$ & $19.135 \pm 5.075$ & $11.834 \pm 3.240^{**}$ & $19.181 \pm 5.096$ \\
Temporal MLP & $20.154 \pm 8.574^{**}$ & $34.996 \pm 15.895$ & $21.119 \pm 8.560^{**}$ & $36.780 \pm 15.935$ & $19.450 \pm 7.934^{**}$ & $33.244 \pm 14.950$ & $19.404 \pm 7.944^{***}$ & $33.138 \pm 14.987$ \\
Naive Persistence & $27.870 \pm 1.299^{***}$ & $37.487 \pm 1.921$ & $28.266 \pm 1.319^{***}$ & $38.202 \pm 1.976$ & $27.697 \pm 1.287^{***}$ & $37.173 \pm 1.881$ & $27.774 \pm 1.289^{***}$ & $37.306 \pm 1.882$ \\
\midrule
\textit{p} vs. ST-GNN (MAE) & \multicolumn{2}{c}{$<0.001$} & \multicolumn{2}{c}{$<0.001$} & \multicolumn{2}{c}{$<0.001$} & \multicolumn{2}{c}{$<0.001$} \\
\bottomrule
\end{tabular}
\end{table*}


Pooled across all four scenarios on the real topology, \CGCEN{} attains
$\mathrm{MAE} = 8.88 \pm 2.40$, against $11.87 \pm 3.12$ for the Spatial Lag
Model, $20.03 \pm 7.97$ for the Temporal MLP, and $27.90 \pm 1.27$ for Naive
Persistence. The RMSE ordering is identical ($13.81$, $19.30$, $34.54$,
$37.54$), indicating that \CGCEN{}'s advantage is not an artefact of a few
well-predicted units but holds under a loss that penalises large residuals more
heavily.

The margin over the graph-blind MLP is substantial and significant in every
scenario, with mean differences between $-10.25$ and $-12.45$ MAE and
$p$-values from $9.4 \times 10^{-4}$ to $5.8 \times 10^{-3}$. The margin over
persistence is larger still ($p < 2 \times 10^{-8}$ throughout). We regard the
persistence comparison as a sanity floor rather than a result: in an earlier
iteration of this work both learned baselines lost to it, which is what first
revealed that the benchmark was mis-specified.

\subsection{Where the Graph Earns Its Keep}

The comparison against the Spatial Lag Model is the substantive one, because
the SLM is itself graph-aware and differs from \CGCEN{} principally in
functional form. Here the result is conditional, and we report the condition
rather than the pooled average:

\begin{center}
\begin{tabular}{lrrr}
\toprule
\textbf{Scenario} & $\Delta$\textbf{MAE} & $t$ & $p$ \\
\midrule
Baseline              & $-1.94$ & $-1.68$ & $0.127$ \\
Market-Led            & $-3.34$ & $-4.14$ & $0.0025$ \\
Inclusionary Housing  & $-3.15$ & $-3.96$ & $0.0033$ \\
Land-Value Capture    & $-3.53$ & $-4.50$ & $0.0015$ \\
\bottomrule
\end{tabular}
\end{center}

\CGCEN{} beats the SLM by roughly $3.3$ MAE under every active intervention,
with large paired effect sizes (Cohen's $d_z$ between $-1.25$ and $-1.42$), but
its $1.94$ MAE advantage in the untreated baseline does not reach significance
at $n = 10$ seeds ($p = 0.127$, $d_z = -0.53$).

This boundary is, we believe, the most useful result in the paper. In a static
equilibrium the system's spatial structure is adequately summarised by a linear
autoregressive lag, and the additional capacity of a nonlinear message-passing
model buys nothing that survives seed variance. It is only when policy
interventions inject treatment-dependent, spatially propagating perturbations
that the nonlinear interference channel $\Psi$ has structure to exploit. A
paper reporting only the pooled comparison would have claimed a general
forecasting advantage that the per-scenario decomposition does not support.

\subsection{Topology Ablation}

To establish that \CGCEN{}'s advantage derives from the graph rather than from
architectural capacity, we re-fit it under two counterfactual topologies while
holding everything else fixed: \emph{no edges}, in which the interference term
of \eqref{eq:decomposition} is identically zero, and \emph{shuffled}, in which
the edge index is remapped through a random node permutation, preserving the
degree sequence and edge count while destroying the correspondence between
graph structure and spatial reality.

\begin{center}
\begin{tabular}{lrrl}
\toprule
\textbf{Topology} & \textbf{MAE} & \textbf{Degradation} & \textbf{$p$ vs.\ real} \\
\midrule
Real graph      & $8.88$  & ---     & --- \\
No edges        & $11.77$ & $+2.89$ & $0.005$--$0.008$ (3/4) \\
Shuffled edges  & $13.76$ & $+4.88$ & $0.0001$--$0.018$ (4/4) \\
\bottomrule
\end{tabular}
\end{center}

Removing the graph costs $2.89$ MAE, a $32.6\%$ relative degradation.
Randomising it costs $4.88$ MAE, a $55.0\%$ degradation, significant in all
four scenarios.

The critical observation is that \emph{shuffled is worse than no edges}. If the
model were merely using message passing as a generic smoothing or
regularisation device, a random graph would serve about as well as the true
one, and both would sit between the real graph and the edgeless condition.
Instead, supplying a structurally plausible but factually wrong topology is
more damaging than supplying none at all. This is the signature of a model that
has learned the specific adjacency: it propagates neighbour information with
confidence, and when that information is drawn from the wrong neighbours the
confident propagation actively injects error.
Figure~\ref{fig:ablation} shows the per-scenario breakdown.

For contrast, the same ablation under the static-shock generator of
Section~\ref{subsec:collapse} gave no-edges $14.16$ against real $14.82$---the
edgeless model was \emph{better}---and shuffled indistinguishable from real
($p = 0.70$--$0.93$). The ablation is thus not merely a confirmation of the
final model; it is the diagnostic that distinguished a working benchmark from a
degenerate one.

\subsection{Causal Channel Behaviour}

Figure~\ref{fig:spillover} visualises the estimated direct and interference
effects on the lattice. The ITE surface concentrates on the treated boundary,
as it must: by \eqref{eq:ite}, $\tau_i$ is identically zero at untreated units.
The STE surface spreads outward from the treated set with magnitude decaying in
graph distance, and is exactly zero at units with no treated neighbour. These
are structural guarantees of the decomposition rather than learned behaviour,
and we verify both as unit tests together with the additive identity
$\hat{Y} = \Phi + \sum_j W_{ij}\Psi$.

\subsection{Policy Scorecard}

Table~\ref{tab:policy_scorecard} ranks the scenarios by displacement-pressure
mitigation. With the level-aware term of \eqref{eq:dpi} active, the contrasts
are substantively legible: Market-Led Development raises $\DPI$ by
$+0.0172 \pm 0.0007$ relative to baseline, while Inclusionary Housing and
Land-Value Capture lower it by $-0.0131 \pm 0.0006$ and $-0.0125 \pm 0.0005$
respectively. The two mitigating policies are close but consistently ordered
across seeds, with inclusionary housing marginally ahead---unsurprising, since
it acts on both rent and housing stock while land-value capture acts on rent
alone.

Under the growth-rate-only specification these same contrasts were on the order
of $10^{-4}$ and returned to zero within one step of the intervention. The
level-aware term increases the measured effect by more than two orders of
magnitude and, as panel (d) of Figure~\ref{fig:trajectories} shows, sustains it
across the full horizon. We stress that this is a change in what the index
\emph{measures}, not a change in the underlying simulated economy: the panels
are identical, and only the readout differs.

\begin{table}[t]
\centering
\caption{Policy scenario scorecard ranked by displacement-pressure mitigation. Negative values indicate that the intervention reduces the index relative to the untreated baseline panel. Values are means over independent seeds.}
\label{tab:policy_scorecard}
\small
\begin{tabular}{clccc}
\toprule
\textbf{Rank} & \textbf{Scenario} & $\Delta$\textbf{DPI} & $\Delta$\textbf{H+T} & $\Delta$\textbf{Rent/Income} \\
\midrule
1 & Inclusionary Housing & $-0.01305 \pm 0.00055$ & $-0.00005 \pm 0.00000$ & $-4.668e-05$ \\
2 & Land-Value Capture & $-0.01255 \pm 0.00054$ & $-0.00004 \pm 0.00000$ & $-4.481e-05$ \\
3 & Baseline & $0.00000 \pm 0.00000$ & $0.00000 \pm 0.00000$ & $0.000e+00$ \\
4 & Market-Led Development & $0.01721 \pm 0.00073$ & $0.00006 \pm 0.00000$ & $6.145e-05$ \\
\bottomrule
\end{tabular}
\end{table}


\begin{figure*}[t]
  \centering
  \includegraphics[width=\textwidth]{spillover_heatmaps.pdf}
  \caption{Direct policy impact versus multi-hop neighbour spillover on the
  $4 \times 4$ lattice under the Market-Led scenario. Top row: treated boundary
  and the $|\ITE|$ surface from $\Phi$. Bottom row: the $|\STE|$ surface from
  the neighbour-aggregated $\Psi$ at three evaluation steps. The ITE is exactly
  zero at untreated units and the STE is exactly zero at units with no treated
  neighbour, both by construction of \eqref{eq:ite}--\eqref{eq:ste}.}
  \label{fig:spillover}
\end{figure*}

\begin{figure*}[t]
  \centering
  \includegraphics[width=\textwidth]{model_comparison_bar.pdf}
  \caption{Model family comparison pooled over ten seeds and four scenarios
  (left, centre) and policy $\DPI$ mitigation relative to the untreated
  baseline (right). Error bars are one standard deviation across seeds.}
  \label{fig:model-comparison}
\end{figure*}

\begin{figure*}[t]
  \centering
  \includegraphics[width=\textwidth]{ablation_topology.pdf}
  \caption{Graph topology ablation. Left: pooled \CGCEN{} accuracy under the
  real graph, an edgeless graph and a randomly rewired graph of identical
  degree sequence. Right: per-scenario breakdown. Randomising the topology is
  more damaging than removing it, indicating the model has learned the specific
  adjacency rather than using message passing as generic smoothing.}
  \label{fig:ablation}
\end{figure*}

\section{Discussion, Limitations and Future Work}
\label{sec:discussion}

\subsection{Interpretation}

Three findings carry beyond the specific model.

First, \textbf{graph structure is not automatically informative}, even when the
data-generating process contains explicit spillover. The static intercept
collapse of Section~\ref{subsec:collapse} shows that spillover which is
permanent and early is absorbed into node-level state and becomes
observationally equivalent to a fixed effect. A benchmark constructed that way
will report that graph models do not help, and the report will be correct about
the benchmark while saying nothing about the method. Practitioners evaluating
spatio-temporal architectures should verify that their spillover process
remains dynamic \emph{within the evaluation window}, and we would encourage
treating the shuffled-edge ablation as a standard diagnostic rather than an
optional extra: it was the shuffled condition, not the edgeless one, that most
clearly separated the working benchmark from the degenerate one.

Second, \textbf{the advantage of nonlinear message passing is
regime-dependent}. In the untreated equilibrium a linear spatial lag is
statistically indistinguishable from \CGCEN{} ($p = 0.127$). Under active
intervention the gap opens to roughly $3.3$ MAE with large effect sizes. The
practical reading is that a planning agency running a steady-state monitoring
model has little reason to prefer a graph neural network over classical spatial
econometrics~\cite{anselin1988,lesage2009}, whereas one comparing prospective
interventions with spatially propagating consequences does.

Third, \textbf{index design can silently determine conclusions}. The
growth-rate Displacement Pressure Index is internally coherent and, applied to
a permanent policy shift, reports approximately zero effect. An analyst who did
not interrogate the functional form could reasonably conclude that inclusionary
housing has no displacement consequence, when in fact the index is structurally
incapable of registering a sustained level difference. The level-aware
correction of \eqref{eq:burden} changes the measured contrast by two orders of
magnitude without altering a single simulated quantity.

\subsection{Limitations}

\subsubsection{No full-scale empirical evaluation}

This is the principal limitation and we state it plainly: \textbf{all
quantitative results in this paper are derived from synthetic data.} We
implement a complete empirical pipeline---a loader that ingests a census-tract
table with ACS and PLUTO fields (median gross rent B25064, median household
income B19013, population B01003, housing units B25001, assessed land value,
lot and residential area, and a derived transit accessibility score), builds
tract adjacency from centroid geodesics by $k$-nearest neighbours under a
haversine metric, and emits a PyTorch Geometric graph---and we validate it end
to end. But we validate it on a \emph{schema-compatible sample}, not on a real
extract. The loader reports its provenance through an explicit
\texttt{is\_synthetic\_sample} flag precisely so that this distinction cannot
be lost downstream, and we have not presented sample-derived numbers anywhere
in this manuscript.

Consequently the results establish \emph{internal} validity---given a
data-generating process with known dynamic spillover, the architecture recovers
it, and we can say exactly when it fails to. They establish nothing about
external validity. Real tract panels are irregular, sparsely observed, subject
to boundary changes across decennial revisions, and contain confounding that
our generator does not simulate. Whether the conditional advantage we document
survives contact with such data is untested, and we would regard any policy
application of \CGCEN{} prior to that test as premature.

\subsubsection{Scale and structure of the synthetic city}

The $4 \times 4$ lattice with $16$ units and $18$ steps yields nine training
windows. This is small, and it interacts with the results in a specific way:
model capacity is cheap to waste, which is part of why the Temporal MLP is so
variable ($\pm 7.97$ MAE) and why the linear SLM is competitive at all.
Real metropolitan graphs are larger, irregular, and have heavy-tailed degree
distributions that a regular lattice does not exercise. We do not claim the
observed ordering transfers unchanged to that setting.

\subsubsection{Simplified interventions and known ground truth}

Policies are applied as multiplicative adjustments on fixed node sets with
fixed timing. Real interventions are heterogeneous in intensity, phased,
politically endogenous to the neighbourhoods that receive them, and often
anticipated by markets before enactment. Treatment endogeneity in particular
would break the exogeneity our design assumes, and addressing it requires
identification machinery---instruments, difference-in-differences with spatial
lags, or an explicit assignment model---that we have not incorporated.

\subsubsection{Partial interference and the exposure summary}

We assume interference operates only through the degree-normalised
one-hop exposure of \eqref{eq:exposure}. This rules out longer-range effects
that are not mediated by intervening units, and it compresses the neighbourhood
assignment vector into a scalar. Both are standard simplifications in the
interference literature~\cite{hudgens2008,forastiere2021}, and both are
falsifiable: richer exposure mappings would be a direct extension.

\subsection{Future Work}

The immediate priority is \textbf{empirical replication}. Sourcing a genuine
NYC ACS/PLUTO tract extract and re-running the benchmark against it is the
single highest-value next step, and the loader exists to make it a
configuration change rather than an engineering project. A useful intermediate
would be a semi-synthetic design: real tract geometry, adjacency and initial
conditions, with simulated interventions layered on, which preserves ground
truth while introducing realistic spatial structure.

Beyond that, three directions follow naturally. \emph{Richer exposure
mappings} would relax the one-hop scalar assumption toward multi-hop or
distance-weighted summaries. \emph{Uncertainty quantification} over the
estimated ITE and STE would let the scorecard report credible intervals on
policy contrasts rather than point estimates; conformal or ensemble approaches
are the obvious candidates. \emph{Treatment assignment modelling} would move
the framework from the ex-ante comparative setting, where assignment is
stipulated, toward observational evaluation, where it must be estimated.

\subsection{Conclusion}

We presented CivicTwin, a counterfactual graph learning framework that
decomposes urban policy effects into a direct term and a spatially
interfering term under partial interference, with the split enforced
architecturally rather than through regularisation. On a controlled synthetic
benchmark across ten seeds and four policy scenarios, \CGCEN{} outperforms a
spatial lag model, a graph-blind temporal MLP and naive persistence, and
topology ablations confirm the advantage is carried by the graph: randomly
rewiring the adjacency is more damaging than removing it entirely.

The advantage is conditional, and we have tried to be precise about the
condition. \CGCEN{} separates from classical spatial econometrics under active
policy intervention ($p = 0.0015$--$0.0033$) but not in the untreated
equilibrium ($p = 0.127$). Under an earlier generator in which spillover
arrived as a single permanent shock, it separated from nothing at all---no
model benefited from the graph, because permanent early spillover collapses
into a static intercept. We report that negative result alongside the positive
one because the boundary between the two regimes is, for anyone building
spatio-temporal policy models, more actionable than either result alone.


\balance
\bibliographystyle{IEEEtran}
\bibliography{paper/references}

\end{document}
